\subsection{Simulation of Unperturbed System}

A Continuous Stirred-Tank Reactor (CSTR) is a chemical engineering system in
which reactants are continuously fed into a well-mixed vessel and products are
continuously withdrawn, maintaining a homogeneous composition throughout the
reactor volume. The CSTR is a standard benchmark for nonlinear control because
its dynamics exhibit strong nonlinearity and coupling between states.

The operation of the discrete-time controlled nonlinear system described by
\eqref{systemModel} and control equation \eqref{eq:control_function}
using a nonlinear cooling CSTR adopted from
\cite{yan2011model,soukkou2008optimal,juang2007recurrent} was
simulated in \cite{tymoshchuk2024tracking}. The dynamics of the reactor can
be described in normalized form by \eqref{systemModel}, where
\begin{equation}
	\label{eq:cstr_dynamics}
	\mathbf{f} =
	\begin{bmatrix}
		x_1(k) + \tau \left( -\alpha x_1(k) + D\bigl(1 - x_1(k)\bigr)
		e^{\dfrac{x_2(k)}{1 + \frac{x_2(k)}{\gamma}}} \right) \\[12pt]
		x_2(k) + \tau \left( -\alpha x_2(k) + BD\bigl(1 - x_1(k)\bigr)
		e^{\dfrac{x_2(k)}{1 + \frac{x_2(k)}{\gamma}}} - \lambda x_2(k) \right)
	\end{bmatrix}, \quad
	\mathbf{g} =
	\begin{bmatrix}
		0 \\ \lambda
	\end{bmatrix},
\end{equation}
with $\alpha = 1.0$, $\lambda = 0.3$, $\gamma = 20.0$,
$B = 1.0$, $D = 0.072$, $\beta = 100$, $\tau = 0.001$\,s,
and $x_1(1) = x_2(1) = 0$.
Here, $x_1$ and $x_2$ denote the dimensionless reactant concentration and
reactor temperature, respectively;
$\alpha$ is the dilution rate;
$D$ is the Damk\"ohler number;
$\gamma$ is the dimensionless activation energy;
$B$ is the dimensionless adiabatic temperature rise;
and $\lambda$ is the dimensionless heat transfer coefficient.
The learning rate $\beta$ is defined in~\eqref{eq:control_function},
and $\tau$ is the discretization time step.

The purpose is to track the reference trajectories
\begin{equation}
	\label{eq:cstr_references}
	r_1(t) =
	\begin{cases}
		0.4472, & \text{if } 0 \leq t \leq 15 \\
		0.7646, & \text{if } 15 < t < 30      \\
		0.4472, & \text{if } 30 \leq t < 45
	\end{cases}
	\quad \text{and} \quad
	r_2(t) =
	\begin{cases}
		2.7520, & \text{if } 0 \leq t \leq 15 \\
		4.7052, & \text{if } 15 < t < 30      \\
		2.7520, & \text{if } 30 \leq t < 45.
	\end{cases}
\end{equation}

The trajectories of state variables $x_1(t)$, $x_2(t)$ obtained based on
\eqref{systemModel} and control equation \eqref{eq:control_function}
and the reference trajectories $r_1(t)$, $r_2(t)$ are presented in
Fig.~\ref{fig:unperturbed-state-tracking}.

\begin{figure}[H]
	\centering
	\includegraphics[width=0.75\textwidth]{cxxplot_time_vs_x1_r1}
	\caption{Trajectory of state $x_1(t)$ and reference $r_1(t)$ for the unperturbed CSTR.}
	\label{fig:unperturbed-state-tracking}
	\label{fig:unperturbed-x1-r1}
\end{figure}

\begin{figure}[H]
	\centering
	\includegraphics[width=0.75\textwidth]{cxxplot_time_vs_x2_r2}
	\caption{Trajectory of state $x_2(t)$ and reference $r_2(t)$ for the unperturbed CSTR.}
	\label{fig:unperturbed-x2-r2}
\end{figure}

The control input $u(k)$ produced by the tracking control law
\eqref{eq:control_function} is shown in Fig.~\ref{fig:unperturbed-control-input}.
The initial transient and the step-change responses at $t = 15$\,s and
$t = 30$\,s are visible, along with the high-frequency switching
(chattering) characteristic of the sliding-mode controller.

\begin{figure}[H]
	\centering
	\includegraphics[width=0.75\textwidth]{cxxplot_fig2b}
	\caption{Trajectory of control $u(t)$.}
	\label{fig:unperturbed-control-input}
	\label{fig:control-input-full}
\end{figure}

\begin{figure}[H]
	\centering
	\includegraphics[width=0.75\textwidth]{cxxplot_fig2b_zoomed}
	\caption{Trajectory of control $u(t)$ for $k \in [0, 100]$.}
	\label{fig:control-input-zoomed}
\end{figure}

The tracking errors are presented in Fig.~\ref{fig:unperturbed-tracking-error}.
Both errors converge to zero between reference changes, confirming the
finite-time convergence property of the control law \eqref{eq:error_dynamics}.

\begin{figure}[H]
	\centering
	\includegraphics[width=0.75\textwidth]{cxxplot_fig2c_e1}
	\caption{Trajectory of tracking error $e_1(t)$.}
	\label{fig:unperturbed-tracking-error}
	\label{fig:tracking-error-e1}
\end{figure}

\begin{figure}[H]
	\centering
	\includegraphics[width=0.75\textwidth]{cxxplot_fig2c_e2}
	\caption{Trajectory of tracking error $e_2(t)$.}
	\label{fig:tracking-error-e2}
\end{figure}
